% ============================================================================
% RAW-TRAJECTORY SNIPPETS -- ROOT REPORT (PLACEHOLDER) -- literal before/after.
% Verbatim ROOT (synthesis) text from instance 27 ("conditions for life on Earth"),
% first vs. last rollout of the v1-clean provenance-smoke RL run. Whitespace
% normalized; [...] marks elision; otherwise VERBATIM. Cited spans are the model's
% own <cite id="...">...</cite> tags, rendered in the provenance colour.
% Source: runs/provenance_smoke/credit_metrics_L3_v1clean.jsonl
% (analysis/training_delta_examples.py --dump-snippets 27, root node).
%
% >>> PLACEHOLDER: regenerate from the FINAL SFT+RL checkpoint before submission. <<<
%
% Self-contained: shares \cit / \trajbox with fig-trajectory-snippets.tex via
% \providecommand (either file may be \input alone or both together).
% Needs \usepackage{tikz} + \usetikzlibrary{arrows.meta} (loaded).
% ============================================================================
\definecolor{cI}{HTML}{1F6FB2}\definecolor{cII}{HTML}{1B8A6B}%
\definecolor{cIII}{HTML}{C97A12}\definecolor{cBad}{HTML}{B03A2E}%
\providecommand{\cit}[2]{\textcolor{cIII}{#2}\,\raisebox{0.3ex}{\scriptsize\textcolor{cIII!85!black}{[\texttt{#1}]}}}
\providecommand{\trajbox}[3]{%
  \begin{tikzpicture}
  \node[draw=#1,line width=0.9pt,rounded corners=3pt,fill=#1!4,
        inner sep=6pt,text width=0.93\columnwidth,align=left,font=\small] (b){#3};
  \node[fill=#1,text=white,font=\scriptsize\bfseries,inner sep=2pt,rounded corners=1pt,
        anchor=west] at ([xshift=6pt,yshift=-0.2pt]b.north west){#2};
  \end{tikzpicture}}

\begin{figure}[t]
\centering
\trajbox{cBad}{BEFORE \textnormal{\scriptsize(start of training)}}{%
\textit{Root report -- ``six elements essential for life'':}\\[2pt]
``\textbf{The Six Essential Elements.} [\,\ldots\,] six are often cited as the most essential
chemical building blocks for life on Earth: Hydrogen, Carbon, Nitrogen, Oxygen,
\textcolor{cBad}{Calcium, and Iron}. These elements form the basis of the organic molecules that
sustain life.''\\[3pt]
{\scriptsize\textcolor{cBad}{$\rightarrow$ \textbf{0 citations}; the ungrounded synthesis even names
the \emph{wrong} six elements (Ca, Fe).}}}

\vspace{4pt}

\trajbox{cII}{AFTER \textnormal{\scriptsize(end of training)}}{%
\textit{Root report -- same sub-questions, grounded synthesis:}\\[2pt]
``[\,\ldots\,] the chemical ingredients are a specific set of elements.
\cit{b5eb73d9-21}{The major macromolecules of the cell account for the bulk of life's mass and are
composed almost entirely of six elements (C,H,N,O,P, and S).}''\\[3pt]
``This planetary habitability is often defined by the Goldilocks Region.
\cit{b5eb73d9-38}{The Goldilocks zone is a region of space where conditions are just right for liquid
water to exist.}''\\[3pt]
{\scriptsize\textcolor{cII!75!black}{$\rightarrow$ \textbf{6 grounded claims}; the synthesis now
cites the \emph{correct} set (CHNOPS), each tag tracing to surfaced evidence.}}}

\caption{\textbf{Root-report output before vs.\ after training (placeholder, v1-clean smoke run).}
Same prompt, verbatim \emph{synthesis} text from the first vs.\ last rollout. The ungrounded draft
asserts (and mis-states) facts; the trained model grounds the same claims in retrieved evidence via
its own \texttt{<cite id>} tags (\textcolor{cIII}{orange}) and gets the chemistry right.
\textit{Final version uses the SFT$+$RL checkpoint.}}
\label{fig:trajectory-snippets-root}
\end{figure}
